\documentclass[a4paper,11pt,twoside]{amsart}
\usepackage[margin=1in]{geometry}
\usepackage{mathtools}
\usepackage{amsfonts}
\usepackage{amsthm}
\usepackage{amsmath}
\usepackage{todonotes}
\usepackage{enumerate}
\usepackage{graphicx}
\usepackage{subcaption}
\usepackage{url}
\usepackage{tabularx}

\let\originalleft\left
\let\originalright\right
\let\vec\mathbf
\renewcommand{\left}{\mathopen{}\mathclose\bgroup\originalleft}
\renewcommand{\right}{\aftergroup\egroup\originalright}
\renewcommand{\Re}{\operatorname{Re}}
\renewcommand{\Im}{\operatorname{Im}}
\newcommand{\Log}{\operatorname{Log}}
\newcommand{\Arg}{\operatorname{Arg}}
\newcommand{\C}{\mathbb{C}}
\newcommand{\R}{\mathbb{R}}
\newcommand{\N}{\mathbb{N}}
\newcommand{\Z}{\mathbb{Z}}
\newcommand{\eps}{\varepsilon}
\newcommand{\Hom}{\operatorname{Hom}}
\newcommand{\catname}[1]{{\textbf{#1}}}

\newcommand\numberthis{\addtocounter{equation}{1}\tag{\theequation}}

\newtheorem{theorem}{Theorem}
\newtheorem{lemma}{Lemma}
\theoremstyle{definition}
\newtheorem*{definition}{Definition}

%\setlength{\oddsidemargin}{.5in}
%\setlength{\evensidemargin}{.5in}
%\setlength{\textwidth}{\paperwidth}
%\addtolength{\textwidth}{-1in}
%\setlength\evensidemargin\oddsidemargin

\begin{document}
\title{(Keeping track of) left and right actions}
\author{Alex Dobner}
\address{Department of Mathematics, University of Michigan \\
530 Church St\\
Ann Arbor, MI 48109}
\email{adobner@umich.edu}
\maketitle

Let $G$ be a group. Recall that a left action of $G$ on a set $X$ is a map $\cdot \colon G\times X \to X$ satisfying the conditions
\begin{enumerate}
    \item $e \cdot x = x$ for all $x \in X$
    \item $g\cdot (h \cdot x) = (gh)\cdot x$ for all $g,h \in G, x \in X$.
\end{enumerate}
We call $X$ a (left) $G$-set.

\textbf{Categorical perspective: } Any group $G$ can be thought of as a category $\ast_G$ with a single object and morphisms for each element of the group. A left group action is then the same thing as a functor $F \colon \ast_G \to \catname{Set}.$ In this case the $G$-set is the unique object in \catname{Set} in the image of the functor.
\bigskip

Alternatively, a \emph{right} action of $G$ on $X$ is a map $\cdot \colon X \times G \to X$ satisfying the conditions
\begin{enumerate}
    \item $x \cdot e = x$ for all $x \in X$
    \item $(x \cdot g) \cdot h = x \cdot (gh)$ for all $g,h \in G, x \in X$.
\end{enumerate}


\vspace{1em}
\textbf{Note:} The reason that the categorical perspective for group actions corresponds to a \emph{left} group action (as opposed to a right group action) is because of how we define composition of functions in conventional mathematics (or equivalently, how we define composition of morphisms in \catname{Set}). Indeed, given two functions $f,g$ we define $f \circ g$ by the rule $(f \circ g)(x) = f(g(x)).$ Right actions correspond to \emph{contravariant} functors from $\ast_G$ to $\mathrm{Set}.$

\smallskip

\textbf{Note 2:} For any group $G$, there is an isomorphism of categories between $\ast_G$ and $\ast_G^{op}$. (This is \emph{not} the same as the trivial fact that any category is \emph{anti}-isomorphic to its opposite category.) This isomorphism allows one to convert a covariant functor from $\ast_G$ to $\catname{Set}$ to a contravariant functor from $\ast_G$ to $\catname{Set}$. In other words, every left $G$ action induces a right $G$ action. This is done by letting $x \cdot g \coloneqq g^{-1} \cdot x.$  Of course one can get a left action from a right action as well.

\section*{Lifting group actions to functions}
Suppose $X$ and $Y$ are sets. If either one of them is is a left or right $G$-set, then there is a natural way to make $\Hom(X,Y)$ into a $G$-set. See the following table for how this goes. (Noting that $g$ denotes a group element and $\Psi \colon X \to Y$ denotes a function.)

\smallskip

\begin{center}
\begin{tabular}{ |c|c|c| }
 \hline
 Type of action: & Lifts to ... on $\Hom(X,Y)$ & Via the definition... \\
 \hline
 left action on $Y$ & a left action & $(g \cdot \Psi)(x) \coloneqq g \cdot (\Psi(x))$ \\
 left action on $X$ & a right action & $(\Psi \cdot g)(x) \coloneqq \Psi(g \cdot x)$ \\
 right action on $Y$ & a right action & $(\Psi \cdot g)(x) \coloneqq \Psi(x) \cdot g$ \\
 right action on $X$ & a left action & $(g \cdot \Psi)(x) \coloneqq \Psi(x \cdot g)$ \\
 \hline
\end{tabular}
\end{center}
\smallskip

The first line corresponds to composing the covariant functor $\ast_G \to \catname{Set}$ (i.e. the left action on $Y$) with the covariant functor $\Hom(X, -) \colon \catname{Set} \to \catname{Set}.$ The second line corresponds to composing the covariant functor $\ast_G \to \catname{Set}$ (the left action on $X$) with the contravariant functor $\Hom(-,Y) \colon \catname{Set} \to \catname{Set}.$ Etc.

\bigskip
The above definitions are what one can get without needing to write a `$g^{-1}$' anywhere (as opposed to what we did in Note~2 above). Composing these rows with the observation in Note~2 gives four more rows...

\begin{center}
\begin{tabular}{ |c|c|c| }
 \hline
 Type of action: & Lifts to ... on $\Hom(X,Y)$ & Via the definition... \\
 \hline
 left action on $Y$ & a right action & $(\Psi \cdot g)(x) \coloneqq g^{-1} \cdot (\Psi(x))$ \\
 left action on $X$ & a left action & $(g \cdot \Psi)(x) \coloneqq \Psi(g^{-1} \cdot x)$ \\
 right action on $Y$ & a left action & $(g \cdot \Psi)(x) \coloneqq \Psi(x) \cdot g^{-1}$ \\
 right action on $X$ & a right action & $(\Psi \cdot g)(x) \coloneqq \Psi(x \cdot g^{-1})$ \\
 \hline
\end{tabular}
\end{center}

\subsection*{Example:} Let $G$ be a compact group. Then we can consider the space $L^2(G)$ of  square integrable functions with respect to Haar measure. The left action of $G$ on itself lifts to a left action of $G$ on $L^2(G)$ using the second row of the second table. This is called the \emph{left regular representation}. The right action of $G$ on itself lifts to a left action of $G$ on $L^2(G)$ using the fourth row of the first table. This is called the \emph{right regular representation}.

%\subsection*{Example} Consider the left action of $GL_2(\C)$ on $\C^2$ given by left matrix multiplication. Projectivizing gives a corresponding left action of $PGL_2(\C)$ on $\C P^1.$


\end{document}
